%=====================================================================
% Datos · Variaciones
%=====================================================================
\subsection*{Variaciones}
\addcontentsline{toc}{subsection}{Variaciones}

Además de los estados, los datos incluyen registros que se salen del caso normal en la forma de su información: columnas que quedan nulas según el tipo de fila, textos que no están en español o que llevan emojis, partidas sin reloj o de cuatro jugadores, registros que no se borran aunque dejen de valer. La consulta siguiente los reúne con una etiqueta que dice qué variación representa cada uno:

\consulta{variaciones}
% Archivo generado por bd/exportar_datos.py. No editar a mano.
\begin{center}\footnotesize\setlength{\tabcolsep}{2pt}
\begin{longtable}{|L{130.6pt}|L{86.3pt}|L{94.5pt}|L{142.0pt}|}
\hline
\textbf{\hspace{0pt}Variación} & \textbf{\hspace{0pt}Tabla} & \textbf{\hspace{0pt}Registro} & \textbf{\hspace{0pt}Dato que la muestra} \\ \hline
\endfirsthead
\hline
\textbf{\hspace{0pt}Variación} & \textbf{\hspace{0pt}Tabla} & \textbf{\hspace{0pt}Registro} & \textbf{\hspace{0pt}Dato que la muestra} \\ \hline
\endhead
Invitado: sin correo ni contraseña & Usuario & Invitado\_\allowbreak{}4821 & correo NULL, contraseña NULL \\ \hline
Cuenta eliminada y anonimizada (D-17) & Usuario & anonimo\_\allowbreak{}10 & correo anonimo\_\allowbreak{}10, baja 2026-08-02 13:00 \\ \hline
Cuenta baneada: sanción permanente & Sancion & 4dm1n\_\allowbreak{}oficial & CUENTA PERMANENTE, fin NULL \\ \hline
Interfaz y términos en inglés (D-21) & AceptacionTerminos & sofia\_\allowbreak{}salto & versión 2026.1 en en \\ \hline
Chat con ñ, acentos, emoji e inglés & Mensaje & mensaje 4 & Welcome! We can start whenever you want \\ \hline
Chat con ñ, acentos, emoji e inglés & Mensaje & mensaje 7 & Igualmente, compañera 👍 \\ \hline
Partida contra la IA, sin reloj & Partida & partida 8 & reloj NULL, deshacer usados 1 \\ \hline
Jugadas deshechas que se conservan (D-12) & Jugada & partida 8 & jugadas 5, 6 con deshecha = 1 \\ \hline
Cuatro jugadores: posiciones y salidas & Participacion & luis\_\allowbreak{}quo & posición 1, salida ninguna \\ \hline
Cuatro jugadores: posiciones y salidas & Participacion & marta\_\allowbreak{}muros & posición 2, salida ABANDONO \\ \hline
Cuatro jugadores: posiciones y salidas & Participacion & pedro\_\allowbreak{}peon & posición 4, salida ABANDONO \\ \hline
Cuatro jugadores: posiciones y salidas & Participacion & sofia\_\allowbreak{}salto & posición 3, salida ABANDONO \\ \hline
Derrota por tiempo: reloj en cero & Participacion & pedro\_\allowbreak{}peon & partida 11, reloj 0 ms \\ \hline
Objeto retirado del catálogo & ObjetoCosmetico & PEON\_\allowbreak{}TEMPORADA\_\allowbreak{}1 & activo 0, a la venta 0 \\ \hline
Término prohibido en todos los idiomas & PalabraProhibida & admin & idioma NULL \\ \hline
Sesión que expiró sin cierre & Sesion & sesión 12 & expiró 2026-08-31 16:00, fin NULL \\ \hline
Sesión que expiró sin cierre & Sesion & sesión 14 & expiró 2026-09-01 16:20, fin NULL \\ \hline
\end{longtable}
\end{center}


\captura{05-variaciones}{Captura 5. Las mismas variaciones en la consola. La consola dibuja el emoji del mensaje 7 como un recuadro porque su fuente no lo tiene.}

Lo que cada grupo permite probar:

\begin{reglas}
  \item \textbf{Nulos que dependen del tipo de fila.} El invitado no tiene correo ni contraseña, y la cuenta eliminada conserva su fila con los datos personales anonimizados (\mbox{D-07}, \mbox{D-17}). Las restricciones \at{CHECK} de subtipo los exigen o los prohíben según el tipo de cuenta.
  \item \textbf{Sanciones permanentes.} La sanción permanente no tiene fecha de fin; la temporal debe tenerla (\mbox{CU-44} \mbox{RN-04}), como comprueba la prueba 11 de los casos que la base rechaza.
  \item \textbf{Idioma y codificación.} \at{sofia\_salto} aceptó los términos en inglés y escribe en inglés; el mensaje 7 lleva una eñe y un emoji. La última comprobación del script de datos verifica que esos textos se guardaron sin pérdida (\mbox{D-21}).
  \item \textbf{Partidas fuera de lo común.} La partida contra la IA no tiene reloj y conserva las dos jugadas que se deshicieron; la de cuatro jugadores guarda la posición de cada uno y la forma en que salió; la perdida por tiempo deja el reloj del perdedor en cero.
  \item \textbf{Lo que deja de valer sin borrarse.} El objeto retirado sigue en el catálogo, sin venderse, porque hay quien lo tiene; las sesiones que expiraron sin cerrarse siguen en su tabla, y el término prohibido sin idioma se filtra en todos.
\end{reglas}

Las jugadas de la partida contra la IA muestran cómo se guarda un deshacer: las jugadas 5 y 6 quedan marcadas como deshechas y la partida sigue desde la posición anterior a ellas, así que en la 7 pedro mueve a otra casilla y en la 8 la IA repite su movimiento (\mbox{D-12}). Las jugadas de la IA no tienen usuario.

\consulta{variaia}
% Archivo generado por bd/exportar_datos.py. No editar a mano.
\begin{center}\scriptsize\setlength{\tabcolsep}{2pt}
\begin{longtable}{|L{50.0pt}|L{54.9pt}|L{84.4pt}|L{42.7pt}|L{46.8pt}|L{35.7pt}|L{85.1pt}|L{36.2pt}|}
\hline
\textbf{\hspace{0pt}Número} & \textbf{\hspace{0pt}Jugador} & \textbf{\hspace{0pt}Tipo} & \textbf{\hspace{0pt}Origen} & \textbf{\hspace{0pt}Destino} & \textbf{\hspace{0pt}Surco} & \textbf{\hspace{0pt}Orientación} & \textbf{\hspace{0pt}Deshecha} \\ \hline
\endfirsthead
\hline
\textbf{\hspace{0pt}Número} & \textbf{\hspace{0pt}Jugador} & \textbf{\hspace{0pt}Tipo} & \textbf{\hspace{0pt}Origen} & \textbf{\hspace{0pt}Destino} & \textbf{\hspace{0pt}Surco} & \textbf{\hspace{0pt}Orientación} & \textbf{\hspace{0pt}Deshecha} \\ \hline
\endhead
1 & pedro\_\allowbreak{}peon & MOVIMIENTO & e1 & e2 & \textit{NULL} & \textit{NULL} & 0 \\ \hline
2 & IA & MOVIMIENTO & e9 & e8 & \textit{NULL} & \textit{NULL} & 0 \\ \hline
3 & pedro\_\allowbreak{}peon & MOVIMIENTO & e2 & e3 & \textit{NULL} & \textit{NULL} & 0 \\ \hline
4 & IA & MURO & \textit{NULL} & \textit{NULL} & e4 & HORIZONTAL & 0 \\ \hline
5 & pedro\_\allowbreak{}peon & MOVIMIENTO & e3 & d3 & \textit{NULL} & \textit{NULL} & 1 \\ \hline
6 & IA & MOVIMIENTO & e8 & e7 & \textit{NULL} & \textit{NULL} & 1 \\ \hline
7 & pedro\_\allowbreak{}peon & MOVIMIENTO & e3 & f3 & \textit{NULL} & \textit{NULL} & 0 \\ \hline
8 & IA & MOVIMIENTO & e8 & e7 & \textit{NULL} & \textit{NULL} & 0 \\ \hline
9 & pedro\_\allowbreak{}peon & MOVIMIENTO & f3 & f4 & \textit{NULL} & \textit{NULL} & 0 \\ \hline
\end{longtable}
\end{center}

